% Shared body of the Week 8 handouts.
% \ifsolutions ... \fi blocks appear only in handout-solutions.tex;
% everything else (statements, key formulas) is identical in both versions.
% Exercise numbers and AK labels follow beamer.tex.

\maketitle

% \begin{remark}
% 	\textbf{How to use this handout.}
% 	Exercise and finance-example numbers match the slides of Week 8.
% 	Problems labeled \textbf{AK01, AK02, \ldots} are taken from the AK
% 	\emph{Probability course, Exercises} collection.
% 	\ifsolutions
% 		Every solution starts with the \emph{idea}, then goes through the steps,
% 		and ends with a boxed answer. Gold boxes point out common mistakes and
% 		alternative methods. Try each exercise yourself before reading the
% 		solution.
% 	\else
% 		A version with detailed solutions is also available. The extra
% 		practice problems at the end of each part are for self-study.
% 	\fi
% 	Questions are welcome by email:
% 	\href{mailto:\myemail}{\texttt{\myemail}}.
% \end{remark}

\tableofcontents
\clearpage

%======================================================================
\section{Part 1: Permutations and combinations}
%======================================================================

\subsection*{Key formulas}
\addcontentsline{toc}{subsection}{Key formulas}

\begin{key}{Counting toolbox}
	\textbf{Basic principle of counting.} If experiment~1 has $m$ outcomes and,
	for each of them, experiment~2 has $n$ outcomes, then together they have
	$m\cdot n$ outcomes. For $r$ experiments: $n_1n_2\cdots n_r$.

	\textbf{Permutations (order matters).} $n$ distinct objects in a row:
	$n! = n(n-1)\cdots 2\cdot1$ (and $0!=1$). Ordered selection of $k$ out of
	$n$ (no repetition):
	\[ \frac{n!}{(n-k)!}=n(n-1)\cdots(n-k+1). \]

	\textbf{Combinations (order does not matter).} Selection of $k$ out of $n$
	as a set (no repetition):
	\[ \binom nk=\frac{n!}{k!\,(n-k)!}=\frac{n!/(n-k)!}{k!}: \]
	each set comes in $k!$ orders, so it was counted $k!$ times among the
	ordered selections.

	\textbf{Also used:} permutations with repetition:
	$n$ objects, $n_i$ identical of kind $i$ ($n_1+\dots+n_r=n$), give
	\[ \frac{n!}{n_1!\,n_2!\cdots n_r!} \]
	different rows.

	% \textbf{Two questions to ask every time.} Does the order matter? Can an
	% element be used more than once?
	% \begin{center}
	% 	\begin{tabular}{lcc}
	% 		\toprule
	% 		                   & order matters & order does not matter       \\\midrule
	% 		no repetition      & $n!/(n-k)!$   & $\binom nk$                 \\
	% 		repetition allowed & $n^k$         & $\binom{n+k-1}{k}$ (Part 2) \\\bottomrule
	% 	\end{tabular}
	% \end{center}

	\textbf{Tricks.} ``Must be together'' $\to$ glue into a block, permute the
	blocks, then permute inside each block. ``Not together'' or ``at least
	one'' $\to$ complement: $\#\text{good}=\#\text{all}-\#\text{bad}$.
\end{key}

%----------------------------------------------------------------------
\subsection{Exercise 1.1 and Finance example 1 --- License plates, PINs}

\begin{task}{Exercise 1.1}
	A license plate has 7 places: the first 2 are letters (26 options each), the
	other 5 are digits. (a) How many plates are possible? (b) How many, if no
	letter and no digit may be repeated?
\end{task}
\workspace{3cm}

\ifsolutions
	\soltitle
	\idea Fill the 7 places one after the other and multiply the numbers of
	choices (basic principle of counting). The order matters (AB12345 and
	BA12345 are different plates); the only question is whether symbols may be
	repeated.

	\step{(a) Repetition allowed.} Each letter place has 26 choices regardless of
	what was chosen before, and each digit place has 10 choices:
	\[ \underbrace{26\cdot26}_{\text{letters}}\cdot
		\underbrace{10\cdot10\cdot10\cdot10\cdot10}_{\text{digits}}
		= 676\cdot100\,000 = 67\,600\,000. \]

	\step{(b) No repetition.} The first letter has 26 choices, the second only
	25 (it must differ from the first). The first digit has 10 choices, the next
	9, then 8, 7, 6:
	\[ (26\cdot25)\cdot(10\cdot9\cdot8\cdot7\cdot6) = 650\cdot30\,240
		= 19\,656\,000. \]
	The letters and the digits are separate pools: a letter never uses up a
	digit, so the digit count starts again from 10.

	\begin{answer}
		(a) $26^2\cdot10^5 = 67\,600\,000$ \qquad
		(b) $26\cdot25\cdot10\cdot9\cdot8\cdot7\cdot6 = 19\,656\,000$
	\end{answer}

	\begin{remark}
		\textbf{Common mistake:} using $\binom{26}{2}$ for the letters. Combinations
		ignore the order, but on a plate the order of the letters matters.
	\end{remark}
\fi

\begin{fin}{Finance example 1 --- PIN codes}
	\context A bank card is protected by a 4-digit PIN; after three wrong
	attempts the card is blocked. The number of possible PINs measures how hard
	it is to guess one.

	(a) How many 4-digit PINs are there? (b) How many have no repeated digit?
	(c) A thief makes three different random guesses. Which fraction of all
	PINs does he cover?
\end{fin}
\workspace{2.5cm}

\ifsolutions
	\soltitle
	\idea A PIN is an ordered string of 4 digits (1234 and 4321 are different
	PINs), so order matters.

	\step{(a)} Each of the 4 places has 10 choices: $10^4=10\,000$ (from 0000 to
	9999).

	\step{(b)} Each digit used removes one option for the next place:
	$10\cdot9\cdot8\cdot7=\dfrac{10!}{6!}=5040$.

	\step{(c)} Three different guesses are 3 of the $10\,000$ equally possible
	PINs, so they cover a fraction $3/10\,000 = 0.0003$. Next week we will call
	this a probability.

	\begin{answer}
		(a) $10\,000$ \qquad (b) $5040$ \qquad (c) $3/10\,000=0.03\%$
	\end{answer}
\fi

%----------------------------------------------------------------------
\subsection{Exercise 1.2 (AK02) --- Five couples at a concert}

\begin{task}{Exercise 1.2 (AK02)}
	Five married couples have 10 seats next to each other in one row. In how
	many ways can they sit (a) with no restriction? (b) if couples sit
	together? (c) if the men sit together and the women sit together? (d) if
	men and women alternate?
\end{task}
\workspace{4cm}

\ifsolutions
	\soltitle
	\idea The 10 people are all different and the seats are in a row, so we count
	permutations. Restrictions of the type ``sit together'' are handled by
	\emph{gluing} people into blocks.

	\step{(a)} Ten different people in ten seats: the first seat has 10 choices,
	the second 9, \ldots, the last 1:
	$10! = 3\,628\,800$.

	\step{(b) Couples together.} Glue each couple into one block. We now arrange
	5 blocks in a row: $5!=120$ orders. Inside each block, the husband can be on
	the left or on the right: 2 orders. This choice is made for each of the 5
	couples independently, giving $2^5=32$. By the counting principle:
	\[ 5!\cdot2^5 = 120\cdot32 = 3840. \]

	\step{(c) Men together, women together.} There are two blocks: the men's
	block and the women's block. They can be ordered in $2$ ways (men on the left
	or women on the left). Inside the men's block the 5 men can be arranged in
	$5!$ ways, and the same for the women:
	\[ 2\cdot5!\cdot5! = 2\cdot120\cdot120 = 28\,800. \]

	\step{(d) Alternating.} First decide the seat pattern: MWMWMWMWMW or
	WMWMWMWMWM, 2 patterns. For a fixed pattern the 5 men fill the 5 ``M'' seats
	in $5!$ ways and the 5 women fill the ``W'' seats in $5!$ ways:
	\[ 2\cdot5!\cdot5! = 28\,800. \]

	\begin{answer}
		(a) $10!=3\,628\,800$ \quad (b) $5!\cdot2^5=3840$ \quad
		(c) $2\cdot5!\cdot5!=28\,800$ \quad (d) $2\cdot5!\cdot5!=28\,800$
	\end{answer}

	\begin{remark}
		(c) and (d) give the same number for a reason: in both cases we first
		choose one of two gender patterns of the seats and then permute the men and
		the women within their seats. \textbf{Common mistake in (b):} forgetting
		the factor $2^5$ for the order inside the blocks.
	\end{remark}
\fi

%----------------------------------------------------------------------
\subsection{Exercise 1.3 (AK03) --- Queueing for the bus}

\begin{task}{Exercise 1.3 (AK03)}
	Eight people line up for a bus. In how many ways can they line up (a) with
	no restriction? (b) if 3 given people insist on standing together? (c) if 2
	given people refuse to stand next to each other?
\end{task}
\workspace{3.5cm}

\ifsolutions
	\soltitle
	\idea (b) is a ``together'' condition $\to$ block. (c) is a ``not
	together'' condition, which is awkward to count directly $\to$ count the
	complement and subtract.

	\step{(a)} $8! = 40\,320$.

	\step{(b)} Glue the 3 friends into one block. Then we arrange 6 units (the
	block and the 5 other people): $6!=720$ ways. Inside the block the 3 friends
	can stand in $3!=6$ orders. Total: $3!\cdot6! = 6\cdot720 = 4320$.

	\step{(c)} Count the \emph{bad} lines, where the two people \emph{are}
	adjacent. Glue them into a block: $7$ units, $7!=5040$ orders, and $2!$
	orders inside the block, so $2\cdot5040=10\,080$ bad lines. All lines minus
	bad lines:
	\[ 8!-2!\cdot7! = 40\,320-10\,080 = 30\,240. \]

	\begin{answer}
		(a) $8!=40\,320$ \quad (b) $3!\cdot6!=4320$ \quad
		(c) $8!-2!\cdot7!=30\,240$
	\end{answer}

	\begin{remark}
		\textbf{Check for (c).} The two people occupy 2 of the 8 positions:
		$\binom82=28$ position pairs, of which 7 are adjacent ($\{1,2\},\dots,
			\{7,8\}$). So a quarter of all lines are bad: $40\,320/4=10\,080$. \
		\textbf{Common mistake:} forgetting the $2!$ inside the block.
	\end{remark}
\fi

%----------------------------------------------------------------------
\subsection{Exercise 1.4 --- Permutations with repetition}

\begin{task}{Exercise 1.4}
	How many different letter arrangements can be made from the letters of
	(a) \mbox{FLUKE} \quad (b) \mbox{PROPOSE} \quad (c) \mbox{MISSISSIPPI} \quad (d) \mbox{ARRANGE}?
\end{task}
\workspace{3cm}

\ifsolutions
	\soltitle
	\idea If all $n$ letters were different there would be $n!$ arrangements.
	If a letter occurs $m$ times, swapping its copies gives the same word, so
	each word has been counted $m!$ times; we divide by the factorials of all
	multiplicities.

	\step{(a) FLUKE.} 5 different letters: $5! = 120$.

	\step{(b) PROPOSE.} 7 letters: P twice, O twice, R, S, E once.
	\[ \frac{7!}{2!\,2!} = \frac{5040}{4} = 1260. \]

	\step{(c) MISSISSIPPI.} 11 letters: M once, I four times, S four times,
	P twice ($1+4+4+2=11$, always check this). Each word is counted
	$4!\cdot4!\cdot2! = 24\cdot24\cdot2 = 1152$ times among the $11!$ orders:
	\[ \frac{11!}{1!\,4!\,4!\,2!} = \frac{39\,916\,800}{1152} = 34\,650. \]

	\step{(d) ARRANGE.} 7 letters: A twice, R twice, N, G, E once:
	$\dfrac{7!}{2!\,2!}=1260$.

	\begin{answer}
		(a) $120$ \quad (b) $1260$ \quad (c) $34\,650$ \quad (d) $1260$
	\end{answer}

	\begin{remark}
		\textbf{Common mistake:} dividing by $2!+2!$ instead of $2!\cdot2!$. The
		over-counting factors multiply: the two P's and the two O's can be swapped
		independently.
	\end{remark}
\fi

%----------------------------------------------------------------------
\subsection{Exercise 1.5 (AK01) and Finance example 2}

\begin{task}{Exercise 1.5 (AK01)}
	The digits $0,1,1,2,3,3,3,4$ are written on eight cards. Using all eight
	cards, how many 8-digit numbers can be formed?
\end{task}
\workspace{3cm}

\ifsolutions
	\soltitle
	\idea Permutations with repetition, minus the arrangements that start with
	$0$ (such a string is not an 8-digit number).

	\step{Step 1: all arrangements.} 8 cards; the digit 1 appears twice and the
	digit 3 three times:
	\[ \frac{8!}{2!\,3!} = \frac{40\,320}{12} = 3360. \]

	\step{Step 2: bad arrangements.} Fix the $0$ in the first place and arrange
	the other 7 cards $1,1,2,3,3,3,4$:
	\[ \frac{7!}{2!\,3!} = \frac{5040}{12} = 420. \]

	\step{Step 3: subtract.}
	\begin{answer}
		$\displaystyle\frac{8!}{2!\,3!}-\frac{7!}{2!\,3!} = 3360-420 = 2940.$
	\end{answer}

	\begin{remark}
		\textbf{Alternative:} by symmetry each of the 8 cards is equally often in
		front, so the $0$ is in front in $1/8$ of the arrangements:
		$3360\cdot\frac78 = 2940$.
	\end{remark}
\fi

\begin{fin}{Finance example 2 --- Up and down days}
	\context Traders often summarize a stock by its daily moves: an
	\emph{up-day} (price closes higher) or a \emph{down-day}. The order of the
	moves describes the price path.

	A trader had 7 up-days and 5 down-days in a 12-day period. How many
	different up/down sequences are possible?
\end{fin}
\workspace{2cm}

\ifsolutions
	\soltitle
	\idea A sequence is a 12-letter word with 7 U's and 5 D's: a permutation
	with repetition. Equivalently: choose which days were down.

	\step{Permutation view.} $\dfrac{12!}{7!\,5!}$.
	\step{Combination view.} Choose the 5 down-days among the 12 days:
	$\binom{12}{5}=\dfrac{12\cdot11\cdot10\cdot9\cdot8}{5!}=\dfrac{95\,040}{120}=792$.
	The two expressions are the same number.
	\begin{answer}
		$\binom{12}{5}=\dfrac{12!}{7!\,5!}=792$
	\end{answer}
\fi

%----------------------------------------------------------------------
\subsection{Exercise 1.6 and Finance example 3 --- Awards and podiums}

\begin{task}{Exercise 1.6}
	Five different awards are given to students of a class of 30. How many
	outcomes are possible if (a) a student may receive any number of awards?
	(b) each student receives at most one award?
\end{task}
\workspace{2.5cm}

\ifsolutions
	\soltitle
	\idea Hand out the awards one by one. The awards are different, so it
	matters who gets which award: order matters. The two parts differ only in
	whether a student can be chosen again.

	\step{(a)} Each of the 5 awards can go to any of the 30 students:
	$30^5 = 24\,300\,000$.

	\step{(b)} The first award has 30 possible recipients, the second 29 (the
	first winner is excluded), then 28, 27, 26:
	\[ 30\cdot29\cdot28\cdot27\cdot26 = \frac{30!}{25!} = 17\,100\,720. \]
	\begin{answer}
		(a) $30^5=24\,300\,000$ \qquad (b) $30!/25!=17\,100\,720$
	\end{answer}
	\begin{remark}
		\textbf{Common mistake:} $5^{30}$ in (a). Ask ``who chooses?'': each
		\emph{award} chooses a student (30 options, 5 times), not each student an
		award.
	\end{remark}
\fi

\begin{fin}{Finance example 3 --- Ranking mutual funds}
	\context Rating agencies and financial newspapers rank investment funds by
	performance; being on the ``podium'' brings new investors to a fund.

	10 mutual funds compete. How many possible results are there if a rating
	agency publishes (a) the gold--silver--bronze podium; (b) only the
	unordered ``top three''; (c) the full league table (ranking of all 10)?
\end{fin}
\workspace{2cm}

\ifsolutions
	\soltitle
	\idea Ask the two questions. A podium is ordered (gold $\ne$ silver) and no
	fund can win two medals; a ``top three'' list is a set; a league table orders
	all 10 funds.
	\step{(a)} Gold: 10 possible funds; silver: 9 (any fund except the gold
	medallist); bronze: 8. Order matters, no fund twice:
	$10\cdot9\cdot8=\dfrac{10!}{7!}=720$.
	\step{(b)} Each unordered top three, e.g.\ $\{X,Y,Z\}$, appears among the
	podiums once for every order of its three members, i.e.\ $3!=6$ times
	($XYZ, XZY, YXZ, YZX, ZXY, ZYX$). Hence
	$720/6=\binom{10}{3}=120$.
	\step{(c)} All 10 funds in order: 10 choices for first place, 9 for second,
	\ldots, 1 for last: $10! = 3\,628\,800$.
	\begin{answer}
		(a) $720$ \quad (b) $120$ \quad (c) $3\,628\,800$
	\end{answer}
\fi

%----------------------------------------------------------------------
\subsection{Exercise 1.7 --- Committees with a conflict}

\begin{task}{Exercise 1.7}
	From 8 women and 6 men a committee of 3 men and 3 women is formed. How many
	committees are possible if (a) 2 of the men refuse to serve together?
	(b) 2 of the women refuse to serve together? (c) 1 man and 1 woman refuse
	to serve together?
\end{task}
\workspace{4cm}

\ifsolutions
	\soltitle
	\idea A committee is a set (no order): combinations. Choose the men and the
	women separately and multiply. For each restriction, count the \emph{bad}
	committees and subtract them from the total.

	\step{Without restriction.} Men: $\binom63=20$, women: $\binom83=56$; total
	$20\cdot56 = 1120$.

	\step{(a)} A bad male trio contains both feuding men; its third member is one
	of the other 4 men: $\binom41=4$ bad trios. Good male trios: $20-4=16$. The
	women are unrestricted: $16\cdot56 = 896$.

	\step{(b)} A bad female trio contains both feuding women plus one of the
	other 6 women: $6$ bad trios. Good: $56-6=50$; total $20\cdot50=1000$.

	\step{(c)} Bad committees contain that man \emph{and} that woman. Then 2
	more men are chosen from the other 5, and 2 more women from the other 7:
	$\binom52\binom72 = 10\cdot21 = 210$. Good: $1120-210=910$.

	\begin{answer}
		(a) $896$ \quad (b) $1000$ \quad (c) $910$
	\end{answer}
	\begin{remark}
		In (c) the restriction links the men and the women, so we must subtract
		from the \emph{total} $1120$, not from the male or female count alone.
		Direct check of (c): committees without the man: $\binom53\binom83=560$;
		with the man but without the woman: $\binom52\binom73=350$;
		$560+350=910$.
	\end{remark}
\fi

%----------------------------------------------------------------------
\subsection{Finance example 4 --- Choosing a portfolio}

\begin{fin}{Finance example 4 --- Choosing a portfolio}
	\context A \emph{portfolio} is a collection of investments. Holding several
	stocks instead of one (\emph{diversification}) reduces risk. Here the
	portfolio is equally weighted, so only the \emph{set} of stocks matters.

	An index has 30 stocks, 8 of them technology stocks. An analyst picks 5
	stocks. (a) How many portfolios are possible? (b) How many contain at least
	one tech stock? (c) How many contain exactly two tech stocks?
	(d) \emph{Lottery.} The Hungarian \emph{ötöslottó} draws 5 numbers out of
	90. How many different tickets are there?
\end{fin}
\workspace{4cm}

\ifsolutions
	\soltitle
	{\emergencystretch=1.5em\idea Unordered selections without repetition: binomial coefficients. ``At
	least one'' $\to$ complement ``none''. ``Exactly two'' $\to$ choose the two
	groups separately.\par}

	\step{(a)}
	$\displaystyle\binom{30}{5}=\frac{30\cdot29\cdot28\cdot27\cdot26}{5\cdot4\cdot3\cdot2\cdot1}
		=\frac{17\,100\,720}{120}=142\,506.$

	\step{(b)} The opposite of ``at least one tech stock'' is ``no tech stock'':
	all 5 stocks come from the $30-8=22$ non-tech stocks, $\binom{22}{5}=26\,334$.
	Hence $142\,506-26\,334=116\,172$.

	\step{(c)} Choose 2 of the 8 tech stocks and 3 of the 22 others:
	$\binom82\binom{22}{3}=28\cdot1540=43\,120$.

	\step{(d)} 5 numbers out of 90, order irrelevant:
	$\binom{90}{5}=43\,949\,268$.

	\begin{answer}
		(a) $142\,506$ \quad (b) $116\,172$ \quad (c) $43\,120$ \quad
		(d) $43\,949\,268$
	\end{answer}
	\begin{remark}
		\textbf{Common mistake in (b):} $8\cdot\binom{29}{4}$ (``choose one tech
		stock, then any 4 others''). This counts a portfolio with two tech stocks
		twice, once for each of them. Use the complement.\\
		\textbf{Lottery:} a single ticket wins the jackpot in 1 of about 44 million
		equally likely draws. Only part of the ticket revenue is paid back as
		prizes, so on average a ticket loses money (negative expected gain) --- we
		will compute such numbers in later weeks.
	\end{remark}
\fi

%----------------------------------------------------------------------
\subsection{Exercise 1.8 --- Dance class}

\begin{task}{Exercise 1.8}
	A dance class has 10 women and 12 men. 5 men and 5 women are chosen and then
	paired off. How many results are possible?
\end{task}
\workspace{2.5cm}

\ifsolutions
	\soltitle
	\idea Split the process into independent steps: choose the women, choose
	the men, match them. Choosing is a combination; matching is a permutation.

	\step{Step 1.} Choose 5 of the 10 women: $\binom{10}{5}=252$.
	\step{Step 2.} Choose 5 of the 12 men: $\binom{12}{5}=792$.
	\step{Step 3.} Line up the 5 chosen women. The first has 5 possible
	partners, the second 4, \ldots, the last 1: $5!=120$ matchings.
	\begin{answer}
		$\displaystyle\binom{10}{5}\binom{12}{5}\,5! = 252\cdot792\cdot120 =
			23\,950\,080$
	\end{answer}
	\begin{remark}
		\textbf{Common mistake:} $5!\cdot5!$ for the matching. Arranging both
		the men and the women counts every set of pairs $5!$ times; only the
		partner of each woman matters.
	\end{remark}
\fi

%----------------------------------------------------------------------
\subsection{Exercise 1.9 --- Lattice paths}

\begin{task}{Exercise 1.9}
	\begin{minipage}[t]{.6\linewidth}
		Starting at $A$, at each move you go one step right or one step up, until
		you reach $B$ (4 steps right, 3 steps up).
		(a) How many paths are there from $A$ to $B$?
		(b) How many of them pass through the circled point $(2,1)$?
	\end{minipage}\hfill
	\begin{minipage}[t]{.36\linewidth}
		\vspace{-4mm}\centering
		\begin{tikzpicture}[scale=.6]
			\draw[gray!50] (0,0) grid (4,3);
			\foreach \x in {0,...,4} \foreach \y in {0,...,3}
			\fill[navy] (\x,\y) circle (2pt);
			\node[below left] at (0,0) {$A$};
			\node[above right] at (4,3) {$B$};
			\draw[teal,very thick] (2,1) circle (6pt);
		\end{tikzpicture}
	\end{minipage}
\end{task}
\workspace{3cm}

\ifsolutions
	\soltitle
	\idea Write a path as a word of moves: R (right) and U (up). Every path from
	$A$ to $B$ has exactly 4 R's and 3 U's, i.e.\ 7 moves, and every such word
	is a path. So we count words.

	\step{(a)} A word is fixed once we decide which 3 of the 7 moves are U's:
	$\binom73=35$. (Or: $\frac{7!}{4!\,3!}=35$.)

	\step{(b)} Split each path at $(2,1)$. From $A=(0,0)$ to $(2,1)$: 2 R's and
	1 U, $\binom31=3$ ways. From $(2,1)$ to $B=(4,3)$: 2 R's and 2 U's,
	$\binom42=6$ ways. Any first half combines with any second half:
	$3\cdot6=18$.

	\begin{center}
		\begin{tikzpicture}[scale=.6]
			\draw[gray!50] (0,0) grid (4,3);
			\foreach \x in {0,...,4} \foreach \y in {0,...,3}
			\fill[navy] (\x,\y) circle (2pt);
			\node[below left] at (0,0) {$A$};
			\node[above right] at (4,3) {$B$};
			\draw[teal,very thick] (2,1) circle (6pt);
			\draw[gold,ultra thick,->] (0,0)--(1,0)--(1,1)--(2,1)--(3,1)--(3,3)--(4,3);
			\node[right] at (4.6,1.5) {\small the path RURRUUR passes through $(2,1)$};
		\end{tikzpicture}
	\end{center}

	\begin{answer}
		(a) $\binom73=35$ \qquad (b) $\binom31\binom42=18$
	\end{answer}
	\begin{remark}
		So $18/35\approx51\%$ of the paths pass through $(2,1)$. The same splitting
		idea is used for binomial trees (Finance example~7 in Part~2): the number
		of paths through a node is (paths into the node) $\times$ (paths out of
		the node). The numbers $\binom nk$ counted here reappear in Pascal's
		triangle, right after Finance example~7.
	\end{remark}
\fi

%----------------------------------------------------------------------
\subsection{Extra practice for Part 1}

\begin{task}{Extra practice 1.1}
	3 boys and 3 girls sit in a row. In how many ways can they sit if (a) there
	is no restriction; (b) the boys sit together and the girls sit together;
	(c) only the boys must sit together; (d) no two people of the same sex sit
	next to each other?
\end{task}
\workspace{2.5cm}

\ifsolutions
	\soltitle
	\idea Same techniques as Exercise~1.2 (AK02): ``together'' $\to$ glue into
	a block; ``alternate'' $\to$ first fix the pattern of the seats.
	\step{(a)} Six different people in six seats: $6!=720$.
	\step{(b)} Two blocks (boys, girls) can be ordered in $2$ ways (boys on the
	left or girls on the left). Inside the boys' block $3!=6$ orders, inside the
	girls' block also $3!=6$: $2\cdot3!\cdot3!=2\cdot6\cdot6=72$.
	\step{(c)} Glue only the boys into one block; together with the 3 girls we
	arrange 4 units: $4!=24$ orders. Inside the block the boys can stand in
	$3!=6$ orders: $24\cdot6=144$.
	\step{(d)} ``No two of the same sex next to each other'' with 3 boys and 3
	girls means that the sexes alternate: BGBGBG or GBGBGB (2 patterns). For a
	fixed pattern the boys fill their 3 seats in $3!$ ways and the girls theirs
	in $3!$ ways: $2\cdot3!\cdot3!=72$.
	\begin{answer}
		(a) $720$ \quad (b) $72$ \quad (c) $144$ \quad (d) $72$
	\end{answer}
	\begin{remark}
		In (c) the girls need not sit together: e.g.\ G\,[BBB]\,G\,G is allowed
		in (c) but not in (b). That is why (c) is larger than (b).
	\end{remark}
\fi

\begin{task}{Extra practice 1.2}
	8 people sit in a row. In how many ways if (a) two given people $A$ and
	$B$ must sit next to each other; (b) there are 4 men and 4 women and no two
	men and no two women sit next to each other; (c) there are 4 married
	couples and each couple must sit together?
\end{task}
\workspace{2.5cm}

\ifsolutions
	\soltitle
	\idea Blocks for ``together'', seat patterns for ``alternating''.
	\step{(a)} Glue $A$ and $B$ into one block. The block and the 6 other people
	are 7 units: $7!=5040$ orders. Inside the block $AB$ or $BA$: 2 orders.
	Total $2\cdot7!=10\,080$.
	\step{(b)} No two men and no two women adjacent means the sexes alternate,
	starting either with a man (\texttt{MWMWMWMW}) or with a woman
	(\texttt{WMWMWMWM}): 2 patterns. The 4 men fill their seats in $4!=24$
	ways, the 4 women theirs in $4!=24$ ways: $2\cdot4!\cdot4!=2\cdot24\cdot24=1152$.
	\step{(c)} Each couple is a block: 4 blocks in $4!=24$ orders. Inside each
	block 2 orders, independently for the 4 couples: $2^4=16$. Total
	$4!\cdot2^4=24\cdot16=384$.
	\begin{answer}
		(a) $10\,080$ \quad (b) $1152$ \quad (c) $384$
	\end{answer}
\fi

\begin{task}{Extra practice 1.3}
	A student has 6 mathematics, 7 science and 4 economics books (all
	different) and sells 2 of them. How many choices are there if the two books
	are on (a) the same subject; (b) different subjects?
\end{task}
\workspace{2.5cm}

\ifsolutions
	\soltitle
	\idea The two books form an unordered pair. Split by subjects and add the
	disjoint cases.
	\step{(a)} Both math: $\binom62=15$; both science: $\binom72=21$; both
	economics: $\binom42=6$; total $15+21+6=42$.
	\step{(b)} One book from each of two different subjects:
	math--science $6\cdot7=42$, math--economics $6\cdot4=24$,
	science--economics $7\cdot4=28$; total $94$.
	\begin{answer}
		(a) $42$ \qquad (b) $94$
	\end{answer}
	\begin{remark}
		Check: every pair is either same-subject or different-subject, and
		$42+94=136=\binom{17}{2}$, the number of all pairs from the 17 books.
	\end{remark}
\fi

\begin{task}{Extra practice 1.4}
	You invite 5 of your 8 friends to a party. How many choices are there if
	(a) 2 of the friends are feuding and will not come together; (b) 2 of the
	friends will only come together (both or neither)?
\end{task}
\workspace{2.5cm}

\ifsolutions
	\soltitle
	\idea The guest list is a set: combinations. (a) is a ``not together''
	condition $\to$ complement; (b) splits into two disjoint cases.
	\step{(a)} Without restriction: $\binom85=56$ guest lists. A bad list
	contains both feuding friends; the other 3 guests come from the remaining
	6 friends: $\binom63=20$ bad lists. Good: $56-20=36$.
	\step{(b)} Call the two friends $X$ and $Y$. \emph{Case 1: both invited.}
	Then 3 more guests from the other 6: $\binom63=20$. \emph{Case 2: neither
		invited.} All 5 guests from the other 6: $\binom65=6$. The cases are
	disjoint, so we add: $20+6=26$.
	\begin{answer}
		(a) $36$ \qquad (b) $26$
	\end{answer}
\fi

\clearpage
%======================================================================
\section{Part 2: Multinomial coefficients, sample space, events}
%======================================================================

\subsection*{Key formulas}
\addcontentsline{toc}{subsection}{Key formulas}

\begin{key}{Multinomial coefficients, expansions, stars and bars}
	\textbf{Multinomial coefficient.} The number of ways to divide $n$ distinct
	objects into $r$ labeled groups of sizes $n_1,\dots,n_r$
	($n_1+\dots+n_r=n$):
	\[ \binom{n}{n_1,\dots,n_r}=\frac{n!}{n_1!\cdots n_r!}
		=\binom{n}{n_1}\binom{n-n_1}{n_2}\cdots\binom{n_r}{n_r}. \]
	\textbf{Binomial and multinomial theorems.}
	\[ (a+b)^n=\sum_{k=0}^n\binom nk a^kb^{n-k},\qquad
		(x_1+\dots+x_r)^n=\sum_{n_1+\dots+n_r=n}\binom{n}{n_1,\dots,n_r}
		x_1^{n_1}\cdots x_r^{n_r}. \]
	\textbf{Binomial tree} (Finance example~7): $\binom nk$ of the $2^n$ paths
	have $k$ up-moves. Pascal's rule (Exercise~2.5) arranges these numbers in
	Pascal's triangle (after Finance example~7).

	\textbf{Combinations with repetition (stars and bars).} $n$ identical
	objects into $r$ distinct boxes: $\binom{n+r-1}{r-1}$ ways; if every box
	must be nonempty: $\binom{n-1}{r-1}$. Equivalently, choosing $k$ elements
	out of $n$ with repetition, order irrelevant: $\binom{n+k-1}{k}$.
\end{key}

\begin{key}{Sample space, events and De Morgan's laws}
	\textbf{Sample space} $\Om$: the set of all possible outcomes of an
	experiment. An \textbf{event} is a subset $E\subseteq\Om$; it
	\emph{occurs} if the outcome lies in $E$.
	$E\cup F$ = ``$E$ or $F$'', \ $E\cap F=EF$ = ``$E$ and $F$'', \ $E^c$ =
	``not $E$'', \ $E\subseteq F$ = ``$E$ implies $F$''. $E$ and $F$ are
	\textbf{mutually exclusive} if $EF=\emptyset$.
	\begin{center}
		\setlength{\tabcolsep}{6mm}
		\begin{tabular}{ccc}
			\venn{\fill[teal!35] (0,0) circle (.85);}
			          & \venn{\fill[teal!35] (1,0) circle (.85);}
			          & \venn{\fill[teal!35] (0,0) circle (.85); \fill[white] (1,0) circle (.85);}                       \\
			$E$       & $F$                                                                        & $E\setminus F=EF^c$ \\[2mm]
			\venn{\fill[teal!35] (0,0) circle (.85); \fill[teal!35] (1,0) circle (.85);}
			          & \venn{\begin{scope}\clip (0,0) circle (.85);
					            \fill[teal!35] (1,0) circle (.85);\end{scope}}
			          & \venn{\shadeall{teal!35}\fill[white] (0,0) circle (.85);}                                        \\
			$E\cup F$ & $EF$                                                                       & $E^c$
		\end{tabular}
	\end{center}

	\textbf{De Morgan's laws for two events.}
	\[ \boxed{(E\cup F)^c=E^c\cap F^c}\qquad\qquad
		\boxed{(E\cap F)^c=E^c\cup F^c} \]
	In words: ``not ($E$ or $F$)'' = ``neither $E$ nor $F$''; ``not (both $E$
	and $F$)'' = ``at least one of them fails''.
	\begin{center}
		\renewcommand{\arraystretch}{1.2}
		\begin{tabular}{ccccc}
			\venn{\shadeall{teal!35}\fill[white] (0,0) circle (.85);}
			                  & $\cap$ &
			\venn{\shadeall{gold!40}\fill[white] (1,0) circle (.85);}
			                  & $=$    &
			\venn{\shadeall{navy!35}\fill[white] (0,0) circle (.85);
				\fill[white] (1,0) circle (.85);}       \\
			\scriptsize $E^c$ &        & \scriptsize $F^c$ &  &
			\scriptsize outside both $=(E\cup F)^c$                       \\[1mm]
			\venn{\shadeall{teal!35}\fill[white] (0,0) circle (.85);}
			                  & $\cup$ &
			\venn{\shadeall{gold!40}\fill[white] (1,0) circle (.85);}
			                  & $=$    &
			\venn{\shadeall{navy!35}\begin{scope}\clip (0,0) circle (.85);
					\fill[white] (1,0) circle (.85);\end{scope}} \\
			\scriptsize $E^c$ &        & \scriptsize $F^c$ &  &
			\scriptsize all but the lens $=(E\cap F)^c$
		\end{tabular}
	\end{center}
	Shade $E^c$ and $F^c$: shaded \emph{twice} $=$ intersection; shaded
	\emph{at least once} $=$ union.

	\textbf{General form ($n$ events).}
	$\big(\bigcup_{i=1}^nE_i\big)^c=\bigcap_{i=1}^nE_i^c$, \
	$\big(\bigcap_{i=1}^nE_i\big)^c=\bigcup_{i=1}^nE_i^c$.
	Next week AK09(c) uses exactly this: $P(A^cB^c)=1-P(A\cup B)$.

	\textbf{Inclusion--exclusion.} $|E\cup F|=|E|+|F|-|EF|$, and
	\[ |E\cup F\cup G|=|E|+|F|+|G|-|EF|-|EG|-|FG|+|EFG|. \]
	Venn diagrams with three sets are filled \emph{from the inside out}.
\end{key}

%----------------------------------------------------------------------
\subsection{Exercise 2.1 (AK04) --- Bridge}

\begin{task}{Exercise 2.1 (AK04)}
	In bridge, the 52 cards are dealt evenly to 4 players (13 each). How many
	deals are possible? Solve it in two ways (permutations and combinations)
	and show that the two answers agree.
\end{task}
\workspace{3.5cm}

\ifsolutions
	\soltitle
	\idea A deal is a division of 52 distinct cards into 4 labeled groups
	(North, East, South, West) of 13 cards each. The order of cards inside a
	hand does not matter.

	\step{Permutation view.} Shuffle the deck: $52!$ orders. Give cards 1--13 to
	player~1, 14--26 to player~2, and so on. Two orders give the same deal if
	they differ only by rearranging cards inside the four blocks of 13; each
	deal therefore arises $13!\cdot13!\cdot13!\cdot13!$ times:
	\[ \frac{52!}{(13!)^4}. \]

	\step{Combination view.} Player~1 gets 13 of the 52 cards, player~2 gets 13
	of the remaining 39, player~3 gets 13 of the remaining 26, player~4 gets the
	last 13:
	\[ \binom{52}{13}\binom{39}{13}\binom{26}{13}\binom{13}{13}. \]

	\step{They agree.} Write out the binomial coefficients; the factorials cancel
	telescopically:
	\[ \frac{52!}{13!\,\cancel{39!}}\cdot\frac{\cancel{39!}}{13!\,\cancel{26!}}
		\cdot\frac{\cancel{26!}}{13!\,13!}\cdot1=\frac{52!}{(13!)^4}. \]
	\begin{answer}
		$\displaystyle\frac{52!}{(13!)^4}=\binom{52}{13}\binom{39}{13}\binom{26}{13}
			\approx5.36\cdot10^{28}$
	\end{answer}
	\begin{remark}
		This is the multinomial coefficient $\binom{52}{13,13,13,13}$. The two
		methods are different (permutations vs.\ combinations) but they count the
		same set, so they must agree.
	\end{remark}
\fi

%----------------------------------------------------------------------
\subsection{Exercise 2.2 and Finance example 5 --- Dividing into groups}

\begin{task}{Exercise 2.2}
	(a) 12 people are divided into 3 committees of sizes 3, 4 and 5. How many
	divisions are possible?
	(b) 8 new teachers are assigned to 4 schools. How many assignments are
	possible? How many, if each school gets exactly 2 teachers?
\end{task}
\workspace{3cm}

\ifsolutions
	\soltitle
	\idea Labeled groups of given sizes $\to$ multinomial coefficient. If the
	group sizes are \emph{not} fixed, each person simply chooses a group.

	\step{(a)} Fill the committees one after the other:
	\[ \binom{12}{3}\binom94\binom55=220\cdot126\cdot1=27\,720
		=\frac{12!}{3!\,4!\,5!}. \]

	\step{(b) Free assignment.} Each of the 8 teachers is sent to one of 4
	schools, independently: $4^8=65\,536$.

	\step{(b) Exactly 2 per school.} Groups of sizes $2,2,2,2$:
	\[ \binom82\binom62\binom42\binom22=28\cdot15\cdot6\cdot1=2520
		=\frac{8!}{2!\,2!\,2!\,2!}. \]
	\begin{answer}
		(a) $27\,720$ \qquad (b) $4^8=65\,536$; with 2 per school $2520$
	\end{answer}
	\begin{remark}
		The schools are different (labeled groups). If the groups were
		\emph{unlabeled} (e.g.\ 8 people split into 4 unnamed pairs), we would
		divide by $4!$ for the order of the groups: $2520/24=105$.
	\end{remark}
\fi

\begin{fin}{Finance example 5 --- Rating bonds}
	\context A \emph{credit rating} (e.g.\ AAA, A, BBB, \ldots, given by
	agencies such as S\&P, Moody's, Fitch) summarizes how likely a borrower is
	to repay. Here an analyst uses three classes A, B, C (best to worst).

	A credit analyst must rate 10 bonds: exactly 4 get rating A, 3 get
	rating B and 3 get rating C. How many rating assignments are possible?
\end{fin}
\workspace{2cm}

\ifsolutions
	\soltitle
	\idea The 10 bonds are different; we divide them into three labeled groups
	(A, B, C) of sizes 4, 3, 3.
	\[ \binom{10}{4,3,3}=\frac{10!}{4!\,3!\,3!}=\binom{10}{4}\binom63\binom33
		=210\cdot20\cdot1=4200. \]
	\begin{answer}
		$4200$
	\end{answer}
	\begin{remark}
		Groups B and C have the same size, but they are \emph{not} interchangeable
		(rating B is not rating C), so we do not divide by $2!$.
	\end{remark}
\fi

%----------------------------------------------------------------------
\subsection{Exercise 2.3 --- Expansions}

\begin{task}{Exercise 2.3}
	(a) Expand $(3x^2+y)^5$. \
	(b) Find the coefficient of $x_1x_2^2x_3$ in $(x_1+2x_2+3x_3)^4$. \
	(c) What is the sum of all coefficients of the expansion in (b)?
\end{task}
\workspace{3.5cm}

\ifsolutions
	\soltitle
	\idea (a) Binomial theorem with $a=3x^2$, $b=y$; the binomial coefficients
	are $\binom5k$: $1,5,10,10,5,1$. (b) Multinomial theorem: each of the 4
	factors contributes one of the three terms. (c) Substitute $1$ for every
	variable.

	\step{(a)} The term with $k$ copies of $a=3x^2$ is
	$\binom5k(3x^2)^ky^{5-k}=\binom5k3^kx^{2k}y^{5-k}$.
	\begin{center}
		\begin{tabular}{ccccc}
			\toprule $k$ & $\binom5k$ & $3^k$ & coefficient & term        \\\midrule
			5            & 1          & 243   & 243         & $243x^{10}$ \\
			4            & 5          & 81    & 405         & $405x^8y$   \\
			3            & 10         & 27    & 270         & $270x^6y^2$ \\
			2            & 10         & 9     & 90          & $90x^4y^3$  \\
			1            & 5          & 3     & 15          & $15x^2y^4$  \\
			0            & 1          & 1     & 1           & $y^5$       \\\bottomrule
		\end{tabular}
	\end{center}
	\[ (3x^2+y)^5=243x^{10}+405x^8y+270x^6y^2+90x^4y^3+15x^2y^4+y^5. \]

	\step{(b)} To get $x_1x_2^2x_3$ we must pick $x_1$ from one factor, $2x_2$
	from two factors and $3x_3$ from one factor. The number of ways to decide
	which factor gives which term is $\binom{4}{1,2,1}=\frac{4!}{1!\,2!\,1!}=12$.
	Each such choice contributes $1\cdot2^2\cdot3=12$. So the coefficient is
	$12\cdot12=144$.

	\step{(c)} The sum of all coefficients is the value of the polynomial at
	$x_1=x_2=x_3=1$: $(1+2+3)^4=6^4=1296$.
	\begin{answer}
		(a) $243x^{10}+405x^8y+270x^6y^2+90x^4y^3+15x^2y^4+y^5$ \quad
		(b) $144$ \quad (c) $1296$
	\end{answer}
	\begin{remark}
		\textbf{Common mistake in (a):} forgetting to raise the $3$ to the power
		$k$ (writing $5\cdot3x^8y$ instead of $5\cdot81x^8y$). \textbf{In (b):}
		forgetting the coefficients $2^2$ and $3$ of the variables.
	\end{remark}
\fi

%----------------------------------------------------------------------
\subsection{Exercise 2.4 --- Identical objects: stars and bars}

\begin{task}{Exercise 2.4 (part (a): AK05)}
	(a) We roll two identical (indistinguishable) dice. How many different
	outcomes can we observe?
	(b) 8 identical blackboards are divided among 4 schools. How many
	divisions are possible? How many, if each school gets at least one?
\end{task}
\workspace{3cm}

\ifsolutions
	\soltitle
	\idea When objects are identical, only the \emph{numbers} matter: how many
	boards each school gets, or which numbers appear on the dice (but not on
	which die).

	\step{(a) First way.} Because the dice look the same, $(1,2)$ and $(2,1)$
	are one observation. Outcomes with two different numbers: choose 2 of 6
	numbers, $\binom62=15$. Doubles: $6$. Total $15+6=21$.

	\step{(a) Second way.} We choose 2 numbers out of 6, order irrelevant,
	repetition allowed: $\binom{6+2-1}{2}=\binom72=21$.

	\step{(b) Stars and bars.} Write the 8 boards as stars and separate the 4
	schools by 3 bars. For example \texttt{**|***||***} means school~1 gets 2,
	school~2 gets 3, school~3 gets 0, school~4 gets 3. Every division gives
	exactly one such code of $8+3=11$ symbols and vice versa, and a code is
	determined by the positions of the 3 bars among the 11 places:
	\[ \binom{11}{3}=165. \]

	\step{(b) At least one each.} Give every school one board first. The
	remaining $8-4=4$ boards are divided freely: 4 stars and 3 bars,
	$\binom{7}{3}=35$.
	\begin{answer}
		(a) $21$ \qquad (b) $\binom{11}3=165$; at least one each:
		$\binom73=35$
	\end{answer}
	\begin{remark}
		\textbf{Common mistake in (a):} $36/2=18$. Halving is wrong because the
		doubles $(1,1),\dots,(6,6)$ appear only once among the 36 ordered
		outcomes, not twice: $(36-6)/2+6=21$.
	\end{remark}
\fi

%----------------------------------------------------------------------
\subsection{Finance example 6 --- Allocating capital}

\begin{fin}{Finance example 6 --- Allocating capital}
	\context \emph{Capital allocation} means deciding how much money goes into
	each investment. Real opportunities often have a minimum ticket size (a
	fund or a project accepts only investments above a threshold).

	We have 20 thousand euros to invest in 4 opportunities, in whole units of
	1000. The minimal investments are 2, 2, 3 and 4 thousand. How many
	investment strategies are there if (a) we invest in each opportunity;
	(b) we invest in at least 3 of the 4? (All 20 units are invested.)
\end{fin}
\workspace{4cm}

\ifsolutions
	\soltitle
	\idea The units of 1000 euros are identical, the opportunities are
	different: stars and bars. Minimums are handled by paying them first.

	\step{(a)} Pay the minimums: $2+2+3+4=11$ units. The remaining $20-11=9$
	units can be divided freely among the 4 opportunities (some may get nothing
	extra): 9 stars and 3 bars,
	\[ \binom{9+3}{3}=\binom{12}{3}=220. \]

	\step{(b)} ``At least 3'' = ``all 4'' (part (a)) or ``exactly 3''. These
	cases are disjoint, so we add. If exactly one opportunity is left out, we pay
	the minimums of the other three and divide the rest among those three
	(2 bars):
	\begin{center}
		\begin{tabular}{lccc}\toprule
			left out    & minimums paid & left to divide & count               \\\midrule
			1st (min 2) & $2+3+4=9$     & 11 among 3     & $\binom{13}{2}=78$  \\
			2nd (min 2) & $2+3+4=9$     & 11 among 3     & $\binom{13}{2}=78$  \\
			3rd (min 3) & $2+2+4=8$     & 12 among 3     & $\binom{14}{2}=91$  \\
			4th (min 4) & $2+2+3=7$     & 13 among 3     & $\binom{15}{2}=105$ \\\bottomrule
		\end{tabular}
	\end{center}
	Total: $220+78+78+91+105=572$.
	\begin{answer}
		(a) $220$ \qquad (b) $572$
	\end{answer}
	\begin{remark}
		\textbf{Common mistake:} in the ``left out'' cases still paying the left-out
		opportunity's minimum, or dividing among 4 boxes instead of 3. A left-out
		opportunity gets exactly $0$.
	\end{remark}
\fi

%----------------------------------------------------------------------
\subsection{Exercise 2.5 (AK07) --- A combinatorial proof}

\begin{task}{Exercise 2.5 (AK07)}
	Prove Pascal's rule: \ $\displaystyle\binom nk+\binom n{k+1}=\binom{n+1}{k+1}$
	\ ($0\le k<n$), (a) by algebra, (b) by a counting argument.
\end{task}
\workspace{4cm}

\ifsolutions
	\soltitle
	\idea (a) Bring the two fractions to a common denominator. (b) Find a
	counting question whose answer is $\binom{n+1}{k+1}$, and answer it a second
	way, by splitting into two cases.

	\step{(a) Algebra.} Write both terms with the common denominator
	$(k+1)!\,(n-k)!$. Since $(k+1)!=(k+1)\,k!$ and $(n-k)!=(n-k)\,(n-k-1)!$:
	\[
		\frac{n!}{k!\,(n-k)!}+\frac{n!}{(k+1)!\,(n-k-1)!}
		=\frac{n!\,(k+1)}{(k+1)!\,(n-k)!}+\frac{n!\,(n-k)}{(k+1)!\,(n-k)!}
		=\frac{n!\,(n+1)}{(k+1)!\,(n-k)!}.
	\]
	Now $n!\,(n+1)=(n+1)!$ and $n-k=(n+1)-(k+1)$, so this equals
	\[ \frac{(n+1)!}{(k+1)!\,\big((n+1)-(k+1)\big)!}=\binom{n+1}{k+1}. \]

	\step{(b) Counting (the ``apple and pear'' argument).} A bowl contains $n$
	apples and $1$ pear, all different. Count the ways to choose $k+1$ fruits.
	\begin{itemize}[nosep]
		\item Directly: $k+1$ out of $n+1$ fruits: $\binom{n+1}{k+1}$.
		\item By whether the pear is chosen. \emph{Pear not chosen}: all $k+1$
		      fruits are apples, $\binom n{k+1}$ ways. \emph{Pear chosen}: the
		      other $k$ fruits are apples, $\binom nk$ ways. These two cases are
		      disjoint and cover everything, so the total is
		      $\binom nk+\binom n{k+1}$.
	\end{itemize}
	The same set was counted in two ways, so the two expressions are equal.
	\begin{answer}
		$\binom nk+\binom n{k+1}=\binom{n+1}{k+1}$
	\end{answer}
	\begin{remark}
		Example: $\binom41+\binom42=4+6=10=\binom52$. After the binomial tree
		(Finance example~7), this rule builds Pascal's triangle. Counting proofs
		avoid long fraction manipulations; see Extra practice 2.2 (AK08).
	\end{remark}
\fi

%----------------------------------------------------------------------
\subsection{Finance example 7 --- The binomial tree}

\begin{fin}{Finance example 7 --- The binomial tree}
	\context In the \emph{binomial (Cox--Ross--Rubinstein) model} of option
	pricing a stock price is modeled in discrete steps: in each step it is
	multiplied either by an up-factor $u>1$ or by a down-factor $d<1$. An
	option's value depends on the final price, so we need to know how many
	paths lead to each final price.

	A stock moves up (factor $u$) or down (factor $d$) every day. After $n=10$
	days the price is $S_0u^kd^{10-k}$, where $k$ is the number of up-days.
	(a) How many price paths are there? (b) How many paths end with $k=6$
	up-days? (c) How many different final prices are there?
\end{fin}
\workspace{3cm}

\ifsolutions
	\soltitle
	\idea A path is a word of 10 letters U/D. The final price depends only on
	how many U's the word contains, not on where they are.

	\step{(a)} 2 choices on each of the 10 days: $2^{10}=1024$.

	\step{(b)} A path with 6 up-days is a word with 6 U's and 4 D's; choose the
	positions of the U's: $\binom{10}{6}=210$.

	\step{(c)} The product $S_0u^kd^{10-k}$ does not depend on the order of the
	factors, so the final price is determined by $k\in\{0,1,\dots,10\}$: 11
	possible prices (they are different because $u>d$).

	\begin{center}
		\begin{tikzpicture}[scale=.85,every node/.style={font=\scriptsize}]
			\foreach \t in {0,...,4}
			\foreach \k in {0,...,\t} {
					\pgfmathtruncatemacro{\c}{factorial(\t)/(factorial(\k)*factorial(\t-\k))}
					\node[circle,draw=navy,fill=navy!8,inner sep=1.5pt,minimum size=13pt]
					(n\t\k) at (1.5*\t,\k-\t/2) {\c};
				}
			\foreach \t in {0,...,3} {
					\pgfmathtruncatemacro{\tn}{\t+1}
					\foreach \k in {0,...,\t} {
							\pgfmathtruncatemacro{\kn}{\k+1}
							\draw[->,gray] (n\t\k)--(n\tn\kn);
							\draw[->,gray] (n\t\k)--(n\tn\k);
						}
				}
			\node[right,align=left,font=\small] at (6.6,0)
			{number of paths reaching each node;\\each node is the sum of the\\two
				feeding it (Pascal's rule)};
		\end{tikzpicture}
	\end{center}

	\begin{answer}
		(a) $2^{10}=1024$ \quad (b) $\binom{10}{6}=210$ \quad (c) $11$
	\end{answer}
	\begin{remark}
		Grouping the 1024 paths by their final node gives
		$\sum_{k=0}^{10}\binom{10}{k}=2^{10}$. In the option-pricing model each path with $k$ up-moves gets
		weight $p^k(1-p)^{10-k}$, and $\binom{10}{k}$ tells how many such paths
		there are --- this is the binomial distribution we will meet later.
		In the figure every node is the sum of the two nodes feeding it: this is
		Pascal's rule of Exercise~2.5, now read as ``paths into a node''.
	\end{remark}
\fi

%----------------------------------------------------------------------
\ifsolutions\else\clearpage\fi % keep the triangle box on one page
\subsection{Pascal's triangle}

\begin{key}{Pascal's triangle}
	\begin{center}
		\begin{tikzpicture}[every node/.style={font=\small}]
			\foreach \n in {0,...,6} {
					\node[navy!70] at (-4.4,-.62*\n) {\scriptsize row $n=\n$};
					\foreach \k in {0,...,\n} {
							\pgfmathtruncatemacro{\c}{factorial(\n)/(factorial(\k)*factorial(\n-\k))}
							\node[inner sep=1pt,minimum size=14pt] (p\n-\k)
							at ({(\k-\n/2)*.95},{-.62*\n}) {\c};
						}
				}
			\draw[gold,thick] (p4-1) circle (7.5pt);
			\draw[gold,thick] (p4-2) circle (7.5pt);
			\draw[teal,thick,fill=teal!15] (p5-2) circle (7.5pt);
			\node[teal] at (p5-2) {10};
			\draw[gold,thick,->] (p4-1)--(p5-2);
			\draw[gold,thick,->] (p4-2)--(p5-2);
			\node[right,align=left,font=\footnotesize] at (3.6,-2.8)
			{$\binom41+\binom42=\binom52$\\ $4+6=10$};
		\end{tikzpicture}
	\end{center}
	Row $n$, position $k=0,1,\dots,n$ (counted from the left, starting at 0)
	holds $\binom nk$. \textbf{Pascal's rule}, just proved in Exercise~2.5
	(AK07),
	\[ \binom nk+\binom n{k+1}=\binom{n+1}{k+1}, \]
	builds the triangle: every entry is the sum of the two entries just above
	it, e.g.\ $\binom41+\binom42=4+6=10=\binom52$.
	\textbf{Readings of row $n$:}
	(1) paths with $k$ up-moves among $n$ moves: the lattice paths of
	Exercise~1.9 and the binomial tree of Finance example~7;
	(2) the coefficients of $(a+b)^n$, e.g.\
	$(a+b)^5=a^5+5a^4b+10a^3b^2+10a^2b^3+5ab^4+b^5$ (row 5);
	(3) the row sum is $2^n$, and each row is symmetric:
	$\binom nk=\binom n{n-k}$.
\end{key}

%----------------------------------------------------------------------
\subsection{Exercise 2.6 and Finance example 8 --- Writing down the sample space}

\begin{task}{Exercise 2.6}
	A box has 3 marbles: red, green and blue. We draw one, put it back, then
	draw a second one. Describe $\Om$. Repeat the exercise if the first marble
	is not put back.
\end{task}
\workspace{2.5cm}

\ifsolutions
	\soltitle
	\idea Record an outcome as an ordered pair (first color, second color),
	e.g.\ $rg$ = red first, green second.

	\step{With replacement.} Both draws can be any of $r,g,b$:
	\[ \Om=\{rr,rg,rb,gr,gg,gb,br,bg,bb\},\qquad|\Om|=3\cdot3=9. \]
	\step{Without replacement.} The second marble cannot have the color of the
	first one, so $rr,gg,bb$ are impossible:
	\[ \Om=\{rg,rb,gr,gb,br,bg\},\qquad|\Om|=3\cdot2=6. \]
	\begin{answer}
		With replacement: 9 outcomes; without replacement: 6 outcomes (listed above).
	\end{answer}
	\begin{remark}
		$rg$ and $gr$ are different outcomes: the draws happen at different
		times. Writing $\{r,g\}$ for both would lose information that we will need
		for probabilities.
	\end{remark}
\fi

\begin{fin}{Finance example 8 --- Events of a stock path}
	\context A stock moves Up or Down on each of 3 days (a three-step binomial
	tree). Questions like ``did the stock end above its starting price?'' or
	``did it ever fall below it?'' are events --- some options (e.g.\
	\emph{barrier options}) pay depending on such path properties.

	$\Om=\{UUU,UUD,UDU,UDD,DUU,DUD,DDU,DDD\}$. Let $A$ = ``the stock ends above
	the start'', $B$ = ``up on day 1'', $C$ = ``never below the start''. List
	$A$, $B$, $C$, and find $AB$, $A^cB$ and $AB^c$. Which inclusion holds
	between $A$ and $C$?
\end{fin}
\workspace{3cm}

\ifsolutions
	\soltitle
	\idea Track the position (number of ups minus number of downs) after each
	day. For example $UDD$: $+1,0,-1$.

	\step{$A$.} The end position is (ups) $-$ (downs) and must be positive, so
	at least 2 U's:
	\[ A=\{UUU,UUD,UDU,DUU\}. \]

	\step{$B$.} First letter U: $B=\{UUU,UUD,UDU,UDD\}$.

	\step{$C$.} Positions: $UUU$: $1,2,3$; $UUD$: $1,2,1$; $UDU$: $1,0,1$;
	$UDD$: $1,0,-1$ (falls below); every path starting with D is at $-1$ after
	day 1. So $C=\{UUU,UUD,UDU\}$.

	\step{Combinations.} $AB=\{UUU,UUD,UDU\}$; $A^cB=\{UDD\}$ (up first but
	ends below); $AB^c=\{DUU\}$ (down first but ends above). Every element of
	$C$ is in $A$, so $C\subseteq A$; in fact $C=AB$.
	\begin{answer}
		$A=\{UUU,UUD,UDU,DUU\}$, $B=\{UUU,UUD,UDU,UDD\}$, $C=\{UUU,UUD,UDU\}$;
		$AB=C$, $A^cB=\{UDD\}$, $AB^c=\{DUU\}$; $C\subseteq A$.
	\end{answer}
	\begin{remark}
		Why $C\subseteq A$: after 3 steps the position is odd ($\pm1$ or $\pm3$),
		so ``not below 0'' at the end means ``at least $+1$'', i.e.\ ending above
		the start.
	\end{remark}
\fi

%----------------------------------------------------------------------
\subsection{Exercise 2.7 --- Two dice: event algebra}

\begin{task}{Exercise 2.7}
	Two dice are rolled. Let $E$ = ``the sum is odd'', $F$ = ``at least one die
	shows 1'', $G$ = ``the sum is 5''. Describe the events $EF$, $E\cup F$,
	$FG$, $EF^c$ and $EFG$.
\end{task}
\workspace{3.5cm}

\ifsolutions
	\soltitle
	\idea $\Om=\{(i,j):i,j=1,\dots,6\}$, $|\Om|=36$. Find the sizes of $E$ and
	$F$ first, then combine them.

	\step{Basic events.} The sum is odd exactly when one die is even and the
	other odd: $3\cdot3$ (first odd, second even) $+\ 3\cdot3$ (first even,
	second odd) $=18$, so $|E|=18$. $F^c$ = ``no die shows 1'' has $5\cdot5=25$
	outcomes, so $|F|=36-25=11$.

	\step{$EF$.} One die shows 1 (odd), so for an odd sum the other must be even:
	\[ EF=\{(1,2),(1,4),(1,6),(2,1),(4,1),(6,1)\},\quad 6\text{ outcomes}. \]
	\step{$E\cup F$.} $|E\cup F|=|E|+|F|-|EF|=18+11-6=23$: all 18 odd-sum
	outcomes plus the outcomes of $F$ with an even sum: $(1,1),(1,3),(1,5),
		(3,1),(5,1)$.
	\step{$FG$.} $G=\{(1,4),(2,3),(3,2),(4,1)\}$; those containing a 1:
	$FG=\{(1,4),(4,1)\}$.
	\step{$EF^c$.} Odd sum and no 1: $|E|-|EF|=18-6=12$ outcomes. One die
	is in $\{2,4,6\}$, the other in $\{3,5\}$ (odd, but not 1): $3\cdot2\cdot2=12$:
	\[ EF^c=\{(2,3),(2,5),(4,3),(4,5),(6,3),(6,5),(3,2),(5,2),(3,4),(5,4),
		(3,6),(5,6)\}. \]
	\step{$EFG$.} $5$ is odd, so $G\subseteq E$ and $EFG=FG=\{(1,4),(4,1)\}$.
	\begin{answer}
		$|EF|=6$, \ $|E\cup F|=23$, \ $FG=\{(1,4),(4,1)\}$, \ $|EF^c|=12$, \
		$EFG=\{(1,4),(4,1)\}$
	\end{answer}
	\begin{remark}
		\textbf{Common mistake:} $|F|=12$ (``6 outcomes with a 1 on the first die,
		6 on the second''). The outcome $(1,1)$ is counted twice; $6+6-1=11$.
	\end{remark}
\fi

%----------------------------------------------------------------------
\subsection{Finance example 9 --- A credit basket}

\begin{fin}{Finance example 9 --- A credit basket}
	\context A \emph{credit basket} is a group of bonds (or loans). A
	\emph{first-to-default swap} is an insurance-like contract that pays when
	the first bond in the basket defaults; a \emph{second-to-default swap} pays
	when a second default occurs. Pricing them starts with describing these
	events precisely.

	A basket contains 3 bonds. Let $D_i$ = ``bond $i$ defaults within a
	year''; $\Om=\{0,1\}^3$ (1 = default), $|\Om|=8$. Express with
	$D_1,D_2,D_3$ and count: (a) at least one default; (b) no default;
	(c) exactly one default; (d) at least two defaults.
\end{fin}
\workspace{3.5cm}

\ifsolutions
	\soltitle
	\idea Translate words into set operations: ``or'' $\to\cup$, ``and''
	$\to\cap$, ``not'' $\to{}^c$. Then count the 0--1 strings.

	\step{(a)} $D_1\cup D_2\cup D_3$ (the first-to-default swap pays). All
	strings except $000$: $8-1=7$.
	\step{(b)} ``No default'' is the complement of (a); by De Morgan
	$(D_1\cup D_2\cup D_3)^c=D_1^cD_2^cD_3^c$ (all three survive): only $000$.
	\step{(c)} Exactly one bond defaults and the other two survive:
	$D_1D_2^cD_3^c\cup D_1^cD_2D_3^c\cup D_1^cD_2^cD_3$. The three pieces are
	mutually exclusive: $100,010,001$, i.e.\ $\binom31=3$ outcomes.
	\step{(d)} Some pair defaults: $D_1D_2\cup D_1D_3\cup D_2D_3$ (the
	second-to-default swap pays): $110,101,011,111$, i.e.\
	$\binom32+\binom33=4$ outcomes.
	\begin{answer}
		(a) $D_1\cup D_2\cup D_3$, 7 \quad (b) $D_1^cD_2^cD_3^c$, 1 \quad
		(c) 3 \quad (d) $D_1D_2\cup D_1D_3\cup D_2D_3$, 4
	\end{answer}
	\begin{remark}
		Checks: (b)+(c)+(d) $=1+3+4=8=|\Om|$ (every outcome has 0, 1 or at least 2
		defaults), and (a) $=$ (c)+(d) $=7$. \textbf{Common mistake in (d):}
		writing $D_1D_2D_3$ --- that is ``all three default'', only one outcome.
	\end{remark}
\fi

%----------------------------------------------------------------------
\subsection{Exercise 2.8 --- Venn diagrams and inclusion--exclusion}

\begin{task}{Exercise 2.8}
	A town of 100\,000 inhabitants has three newspapers I, II, III. The
	percentages of readers: I: 10, II: 30, III: 5, I and II: 8, I and III: 2,
	II and III: 4, all three: 1. I and III are morning papers, II is an
	evening paper. How many people read (a) exactly one paper? (b) at least two
	papers? (c) at least one morning paper and the evening paper? (d) no paper?
	(e) exactly one morning paper and the evening paper?
\end{task}
\workspace{4cm}

\ifsolutions
	\soltitle
	\idea The given numbers overlap (``I and II'' includes those who also read
	III). Fill the Venn diagram \emph{from the inside out}; then every region
	holds the people of exactly that combination, and we can simply add regions.

	\begin{minipage}{.58\linewidth}
		\step{Filling the diagram (in \%).}
		\begin{itemize}[nosep]
			\item center (all three): $1$;
			\item only I\&II: $8-1=7$; only I\&III: $2-1=1$; only II\&III: $4-1=3$;
			\item only I: $10-7-1-1=1$;\\ only II: $30-7-3-1=19$;\\ only III:
			      $5-1-3-1=0$.
		\end{itemize}
	\end{minipage}\hfill
	\begin{minipage}{.38\linewidth}
		\centering
		\begin{tikzpicture}[scale=.65]
			\draw[navy,thick] (0,0) circle (1.5);
			\draw[teal,thick] (1.8,0) circle (1.5);
			\draw[gold,thick] (.9,-1.5) circle (1.5);
			\node at (-.7,.3) {\small 1};
			\node at (2.5,.3) {\small 19};
			\node at (.9,-2.3) {\small 0};
			\node at (.9,.6) {\small 7};
			\node at (.1,-.95) {\small 1};
			\node at (1.7,-.95) {\small 3};
			\node at (.9,-.45) {\small 1};
			\node[navy] at (-1.4,1.4) {I};
			\node[teal] at (3.2,1.4) {II};
			\node[gold] at (.9,-3.3) {III};
		\end{tikzpicture}
	\end{minipage}

	\step{(a)} Exactly one: $1+19+0=20\%$, i.e.\ $20\,000$ people.
	\step{(b)} At least two: the four inner regions $7+1+3+1=12\%$: $12\,000$.
	\step{(c)} $(\mathrm{I}\cup\mathrm{III})\cap\mathrm{II}=
		(\mathrm{I}\cap\mathrm{II})\cup(\mathrm{III}\cap\mathrm{II})$. Inclusion--exclusion:
	$8+4-1=11\%$ (the readers of all three were counted twice). Or from the
	diagram: $7+3+1=11\%$: $11\,000$.
	\step{(d)} $|\mathrm{I}\cup\mathrm{II}\cup\mathrm{III}|=10+30+5-8-2-4+1=32\%$
	(check: the seven regions sum to $1+19+0+7+1+3+1=32$), so $100-32=68\%$ read
	none: $68\,000$.
	\step{(e)} II together with exactly one of I, III (not both): only I\&II
	plus only II\&III: $7+3=10\%$: $10\,000$.
	\begin{answer}
		(a) $20\,000$ \quad (b) $12\,000$ \quad (c) $11\,000$ \quad
		(d) $68\,000$ \quad (e) $10\,000$
	\end{answer}
	\begin{remark}
		\textbf{Common mistake:} computing ``only I'' as $10-8-2=0$. This
		subtracts the center (readers of all three) twice. Always start from the
		center.
	\end{remark}
\fi

%----------------------------------------------------------------------
\subsection{Finance example 10 --- Bank products}

\begin{fin}{Finance example 10 --- Bank products}
	\context Banks try to sell additional products to existing clients
	(\emph{cross-selling}), because an existing client is cheaper to reach than
	a new one. Knowing who already has which product tells the bank whom to
	target.

	A bank has 1000 retail clients: 820 have a current account ($C$), 350 a
	credit card ($K$), 240 a mortgage ($M$); 300 have both $C$ and $K$, 200
	have $C$ and $M$, 90 have $K$ and $M$, and 70 have all three. (a) How many
	clients have none of the products? (b) How many have exactly one product?
	(c) A campaign targets clients with a current account but no credit card.
	How many are they?
\end{fin}
\workspace{4cm}

\ifsolutions
	\soltitle
	\idea Inclusion--exclusion for three sets, and ``from the inside out'' for the
	``only'' regions.

	\step{(a)} $|C\cup K\cup M|=820+350+240-300-200-90+70=890$. None:
	$1000-890=110$.

	\step{(b)} ``Only $C$'' $=|C|-|CK|-|CM|+|CKM|$: removing $CK$ and $CM$
	removes the 70 clients with all three products twice, so we add them back
	once. Only $C$: $820-300-200+70=390$; only $K$: $350-300-90+70=30$; only
	$M$: $240-200-90+70=20$. Exactly one product: $390+30+20=440$.

	\step{(c)} $C$ splits into $CK$ (has a card) and $CK^c$ (no card):
	$|CK^c|=|C|-|CK|=820-300=520$.
	\begin{answer}
		(a) $110$ \quad (b) $440$ \quad (c) $520$
	\end{answer}
	\begin{remark}
		Dividing by 1000 turns these counts into probabilities for a randomly
		chosen client; e.g.\ $P(C\cup K\cup M)=0.89$. This is exactly how
		probabilities will be computed next week.
	\end{remark}
\fi

%----------------------------------------------------------------------
\subsection{Extra practice for Part 2}

\begin{task}{Extra practice 2.1 (AK06)}
	Using the binomial theorem, show that
	$\displaystyle\sum_{k=0}^n\binom nk=2^n$, \
	$\displaystyle\sum_{k=0}^n\binom nk2^k=3^n$ \ and \
	$\displaystyle\sum_{k=0}^n(-1)^k\binom nk=0$ \ ($n\ge1$).
\end{task}
\workspace{3cm}

\ifsolutions
	\soltitle
	\idea Substitute suitable numbers into
	$(a+b)^n=\sum_{k=0}^n\binom nk a^kb^{n-k}$.
	\step{First identity.} $a=b=1$: every term becomes $\binom nk$, and the left
	side is $2^n$.
	\step{Second.} $a=2$, $b=1$: $\sum_k\binom nk2^k=(2+1)^n=3^n$.
	\step{Third.} $a=-1$, $b=1$: $\sum_k\binom nk(-1)^k=(-1+1)^n=0^n=0$ for
	$n\ge1$.
	\begin{answer}
		$\sum_k\binom nk=2^n$, \ $\sum_k\binom nk2^k=3^n$, \
		$\sum_k(-1)^k\binom nk=0$.
	\end{answer}
	\begin{remark}
		Counting meaning of the first identity: both sides count all subsets of an
		$n$-element set (by size on the left; ``in or out'' for each element on the
		right). The third says: an $n$-element set ($n\ge1$) has as many
		even-sized as odd-sized subsets, e.g.\ for $n=4$: $1+6+1=4+4$.
	\end{remark}
\fi

\begin{task}{Extra practice 2.2 (AK08)}
	Prove Vandermonde's identity
	$\displaystyle\sum_{j=0}^k\binom nj\binom m{k-j}=\binom{n+m}k$
	by a counting argument (apples and pears).
\end{task}
\workspace{3cm}

\ifsolutions
	\soltitle
	\idea Count the same set in two ways, as in Exercise~2.5.
	\step{Setting.} A bowl contains $n$ different apples and $m$ different
	pears. Count the ways to choose $k$ fruits.
	\step{Directly.} $k$ out of $n+m$ fruits: $\binom{n+m}k$.
	\step{By the number of apples.} Let $j$ be the number of apples chosen
	($j=0,1,\dots,k$). Then $k-j$ pears are chosen: $\binom nj\binom m{k-j}$
	ways. The cases $j=0,\dots,k$ are disjoint and cover all choices, so the
	total is $\sum_{j=0}^k\binom nj\binom m{k-j}$.
	\begin{answer}
		Both expressions count the $k$-element selections from $n+m$ fruits,
		hence $\sum_{j=0}^k\binom nj\binom m{k-j}=\binom{n+m}k$.
	\end{answer}
	\begin{remark}
		Terms with $j>n$ or $k-j>m$ are $0$ by convention, so the identity holds
		for all $k$. Check with $n=4$, $m=5$, $k=3$:
		$1\cdot10+4\cdot10+6\cdot5+4\cdot1=84=\binom93$.
		Pascal's rule (Exercise~2.5) is the special case $m=1$.
	\end{remark}
\fi

\begin{task}{Extra practice 2.3}
	$A$, $B$ and $C$ take turns flipping a coin, $A$ first, then $B$, then $C$,
	then $A$ again, and so on. The first one to get a head wins. Write down the
	sample space and the events ``$A$ wins'', ``$B$ wins'' and $(A\cup B)^c$.
\end{task}
\workspace{3cm}

\ifsolutions
	\soltitle
	\idea Record the sequence of flips until the first head: 1 = head, 0 =
	tail. The game stops at the first 1, so an outcome is some tails followed by
	one head --- or tails for ever.
	\step{Sample space.}
	\[ \Om=\{1,\ 01,\ 001,\ 0001,\ 00001,\ \dots\}\cup\{000\cdots\}. \]
	A string of $m$ zeros followed by a 1 means that the first head comes at flip
	$m+1$; the infinite string $000\cdots$ means that no head ever comes.
	\step{Who flips when?} $A$ flips at flips $1,4,7,\dots$, $B$ at
	$2,5,8,\dots$, $C$ at $3,6,9,\dots$.
	\step{Events.} ``$A$ wins'' $=\{1,\ 0001,\ 0000001,\dots\}$ (the number of
	leading zeros is a multiple of 3); ``$B$ wins'' $=\{01,\ 00001,\ 00000001,
		\dots\}$ (number of zeros $\equiv1\pmod 3$).
	$(A\cup B)^c$ = neither $A$ nor $B$ wins = $C$ wins or nobody wins:
	$\{001,\ 000001,\dots\}\cup\{000\cdots\}$.
	\begin{answer}
		$A$ wins: $3j$ zeros then 1; $B$ wins: $3j+1$ zeros then 1;
		$(A\cup B)^c$: $3j+2$ zeros then 1, or $000\cdots$ \ ($j=0,1,2,\dots$).
	\end{answer}
	\begin{remark}
		\textbf{Common mistake:} forgetting the outcome $000\cdots$ (no head ever).
		It is in the sample space; next week we will see that its probability is
		0, but it still belongs to $(A\cup B)^c$.
	\end{remark}
\fi

\begin{task}{Extra practice 2.4}
	For each of the 15 members of a sports team we record the type of job
	(blue-collar or white-collar) and the political affiliation (Republican,
	Democrat or Independent). How many outcomes are (a) in the sample space?
	(b) in the event ``at least one member is blue-collar''? (c) in the event
	``no member is Independent''?
\end{task}
\workspace{3cm}

\ifsolutions
	\soltitle
	\idea Each member gets one of $2\cdot3=6$ classifications, and the outcome
	lists the classification of each of the 15 (distinct) members.
	\step{(a)} $6^{15}=470\,184\,984\,576$.
	\step{(b)} Complement: ``nobody is blue-collar'' means every member is
	white-collar with one of 3 affiliations: $3^{15}=14\,348\,907$. So
	$6^{15}-3^{15}=470\,170\,635\,669$.
	\step{(c)} Each member has $2$ job types and $2$ affiliations (R or D):
	$4^{15}=1\,073\,741\,824$.
	\begin{answer}
		(a) $6^{15}$ \quad (b) $6^{15}-3^{15}$ \quad (c) $4^{15}$
	\end{answer}
	\begin{remark}
		\textbf{Common mistake in (b):} $15\cdot6^{14}$ (``choose the blue-collar
		member, the others are free''): outcomes with several blue-collar members
		are counted many times. ``At least one'' $\to$ complement.
	\end{remark}
\fi

\begin{task}{Extra practice 2.5 (reliability of a system)}
	A system has 5 components, each working (1) or failed (0); the outcome is
	$(x_1,\dots,x_5)$. The system works if components 1 and 2 work, or 3 and 4
	work, or 1, 3 and 5 work. (a) What is $|\Om|$? (b) List $W$ = ``the system
	works''. (c) Let $A$ = ``components 4 and 5 failed''. What is $|A|$?
	(d) List $AW$.

	\context Read ``component working'' as ``issuer has not defaulted''. A
	\emph{structured product} built on issuers 1--5 keeps paying as long as at
	least one route survives, i.e.\ as long as the system works. (e) Does the
	product keep paying if any single issuer defaults? Which pairs of defaults
	stop the payments?
\end{task}
\workspace{4cm}

\ifsolutions
	\soltitle
	\idea The system is three parallel routes; a route works if all its
	components work. To list $W$ without double counting, split it into
	\emph{disjoint} pieces.

	\begin{center}
		\begin{tikzpicture}[x=1cm,y=.8cm,every node/.style={font=\small},
				c/.style={draw,circle,fill=navy!10,inner sep=1.5pt}]
			\coordinate (in) at (0,0); \coordinate (out) at (5,0);
			\node[left] at (in) {in}; \node[right] at (out) {out};
			\foreach \y in {1,0,-1} {
					\draw (in) |- (.6,\y) -- (4.4,\y) -| (out);
				}
			\node[c] at (1.8,1) {1}; \node[c] at (3.2,1) {2};
			\node[c] at (1.8,0) {3}; \node[c] at (3.2,0) {4};
			\node[c] at (1.4,-1) {1}; \node[c] at (2.5,-1) {3};
			\node[c] at (3.6,-1) {5};
		\end{tikzpicture}
	\end{center}

	\step{(a)} Each $x_i\in\{0,1\}$: $|\Om|=2^5=32$.

	\step{(b)} $W=\{x_1=x_2=1\}\cup\{x_3=x_4=1\}\cup\{x_1=x_3=x_5=1\}$.
	\begin{itemize}[nosep]
		\item Piece 1: $x_1=x_2=1$, the others free: $2^3=8$ outcomes:
		      \texttt{11000, 11001, 11010, 11011, 11100, 11101, 11110, 11111}.
		\item Piece 2: $x_3=x_4=1$ but \emph{not} $x_1=x_2=1$: $(x_1,x_2)\in
			      \{00,01,10\}$ (3 options), $x_5$ free (2 options): $6$ outcomes:
		      \texttt{00110, 00111, 01110, 01111, 10110, 10111}.
		\item Piece 3: works only through route $1,3,5$: then $x_1=x_3=x_5=1$,
		      and routes 1--2 and 3--4 are broken, so $x_2=0$ and $x_4=0$:
		      \texttt{10101}.
	\end{itemize}
	$|W|=8+6+1=15$.

	\step{(c)} $A=\{x_4=x_5=0\}$, with $x_1,x_2,x_3$ free: $|A|=2^3=8$.

	\step{(d)} If 4 and 5 failed, routes 3--4 and 1--3--5 are broken, so the
	system works only if $x_1=x_2=1$; $x_3$ is free:
	$AW=\{11000,\,11100\}$.
	\begin{answer}
		(a) $32$ \quad (b) $|W|=15$ (listed above) \quad (c) $|A|=8$ \quad
		(d) $AW=\{11000,11100\}$
	\end{answer}
	\begin{remark}
		Check: $|W^c|=32-15=17$ states in which the system fails.
	\end{remark}

	\step{(e) Finance context.} The product stops paying exactly when
	\emph{every} route contains a defaulted issuer. One default breaks at most the
	routes through that issuer: issuer 1 lies on routes 1--2 and 1--3--5 (route
	3--4 survives), issuer 3 on routes 3--4 and 1--3--5 (route 1--2 survives),
	issuers 2, 4, 5 on one route each. So any single default is survived. A pair
	of defaults must hit route $\{1,2\}$ and route $\{3,4\}$, so it is one of
	$\{1,3\},\{1,4\},\{2,3\},\{2,4\}$; of these, $\{1,3\}$, $\{1,4\}$ and
	$\{2,3\}$ also hit route $\{1,3,5\}$, but $\{2,4\}$ does not.
	\begin{answer}
		(e) Yes, every single default is survived; exactly the pairs $\{1,3\}$,
		$\{1,4\}$, $\{2,3\}$ stop the payments.
	\end{answer}
	\begin{remark}
		Issuers 1 and 3 appear on two routes each: they matter most for the
		product, like critical components in an engineering system.
	\end{remark}
\fi
